\documentclass[11pt]{article}

\usepackage[margin=1in]{geometry}

\setlength{\parindent}{0pt}
\setlength{\parskip}{0.6em}
\setlength{\emergencystretch}{2em}

\begin{document}

\begin{center}
    {\Large \textbf{Meeting Summary}}\\[0.4em]
    {\large 14 August 2026}\\[0.4em]
    Bulun Te and Lifeng Lin
\end{center}

\section*{Meeting Summary}

The meeting focused on four research workstreams: completion of a Bayesian heterogeneity manuscript, final preparation of the TAILOR manuscript, expansion of an AI-assisted audit of reporting errors in medical publications, and possible projects in component network meta-analysis (CNMA). The near-term priorities are to integrate the new Bayesian case study, complete one more audit of TAILOR, and rerun the publication-error workflow with a less restrictive prompt before arranging a joint meeting with Shiying.

\subsection*{Bayesian $I^2$ Manuscript and Case Study}

The newly prepared case study is suitable for the manuscript and should replace one of the three existing examples rather than create a fourth example. It analyzes diarrhea during acetylcholinesterase-inhibitor treatment versus placebo in patients with Alzheimer's disease or Parkinson's disease dementia. The extracted dataset contains 37 trial comparisons, with 989 events among 10,825 treated participants and 397 events among 6,863 placebo participants. The Bayesian analysis estimates $I^2=50.7\%$ (95\% credible interval, $26.3\%$--$70.4\%$), while the frequentist DerSimonian--Laird analysis estimates $I^2=43.3\%$ (95\% confidence interval, $9.9\%$--$60.6\%$). The pooled odds ratios are 1.61 (95\% credible interval, 1.33--1.99) and 1.58 (95\% confidence interval, 1.31--1.90), respectively.

The example was accepted because it provides an intermediate-heterogeneity setting and a larger number of studies than the existing examples. Bulun will merge the separate case-study draft into the main Overleaf manuscript, replace the middle example, and explain the selection of examples more clearly. The new forest plot should be presented cleanly as its own figure, without making an existing multi-panel display unwieldy. The analysis and graphical code should be reconciled with the original shared code to preserve consistency across all examples. Bulun will also read the complete manuscript and perform a final check of terminology, punctuation, and technical consistency. Subject to this revision and a final review, the manuscript can be submitted; \emph{Statistics in Biosciences} was discussed as a prospective outlet, but the journal choice remains tentative.

A package implementation is a later deliverable rather than a prerequisite for submission. The intended package contribution would consist of a simple R function, the three example datasets, and their documentation. A full S3 interface is not required for the initial computational function. Function naming, package-wide cleanup, and the timing of integration remain to be decided after the manuscript work is complete and after Lifeng's planned package update.

\subsection*{TAILOR Manuscript}

The revisions requested at the 24 July meeting have been completed. The manuscript is back in the journal's official single-column format; the Student $t$ prior now explains why the default degrees-of-freedom range is user-modifiable and why its lower bound is slightly above 2; repeated equations and descriptions have been replaced with labels and cross-references; and the Shiny section has been refocused on workflow, diagnostics, and outputs. The application sections now provide clinical background, treatment and comparator definitions, tumor-response outcomes, data sources, and the results of the original meta-analyses.

The revised text also distinguishes the risk ratios reported in the original reviews from the odds ratios estimated by TAILOR, so comparisons are restricted to direction and uncertainty rather than numerical equivalence. It clarifies that the empirical lognormal heterogeneity priors were calibrated for log odds ratios and should not be transferred unchanged to other effect measures. Terminology and typography have also been standardized.

Lifeng agreed that the manuscript is essentially ready for submission. Bulun will run additional AI-assisted checks and inspect any flagged issue manually. After Bulun reports that this audit is complete, Lifeng will perform a final review, prepare the cover letter, and handle the initial submission. Submission in the following week was considered feasible, but no fixed deadline was assigned.

\subsection*{AI-Assisted Audit of Published Papers}

The completed 2025 audit currently covers 14 papers and contains 71 verified findings: 35 presentation inconsistencies, 15 statistical-reporting inconsistencies, 10 arithmetic inconsistencies, 10 cross-document inconsistencies, and one participant-flow inconsistency. Eight additional candidates remain explicitly uncertain and should not be treated as established errors. The reviewed examples included mislabeled effect measures, inconsistent analysis or randomization units, duplicated panels, sign errors, point estimates lying outside their confidence intervals, denominator discrepancies, incorrect totals, and conflicts between main articles and supplements.

The current instruction to return exactly ten issues per paper should be removed. That limit can suppress findings in error-rich papers and encourage weak or artificial findings in papers with few genuine problems. The prompt should instead define the desired problem types carefully and ask for all supportable findings. Numerical and statistically relevant inconsistencies will remain the principal focus, while design or analysis-unit problems may be retained when they have a clear statistical consequence. Ambiguous findings must remain qualified and separated from verified results.

The time comparison should cover balanced pre- and post-periods: 2018--2019 and 2024--2025. This design may describe changes across the pre-COVID, post-COVID, and widely available AI eras, but causal claims about either COVID-19 or AI use are not supported by the meeting evidence. Once all papers have been rerun, the human-verification work will be divided with Shiying so that both time periods receive review. The exact allocation will be decided at the joint meeting.

For reporting, detailed article-level findings and model reasoning should be organized by year and paper in supplementary material. The main manuscript should summarize the distribution of issue types and use a small number of well-supported examples to demonstrate scientific or interpretive consequences. The intended tone is constructive: even minor inconsistencies matter because they can enter later reviews and evidence syntheses, and a lightweight AI-assisted check may prevent them. The project should not imply that every flagged issue changes a paper's conclusions.

\subsection*{Potential CNMA Projects and Research Direction}

Two CNMA projects were considered. The first would reanalyze a recent JAMA network meta-analysis of multicomponent interventions using CNMA and compare the resulting conclusions with those from traditional network meta-analysis. Possible comparisons include model fit, such as the deviance information criterion in a Bayesian analysis, or an appropriate frequentist component model. This project is conditional on obtaining analysis-ready data. Re-extracting a large dataset from publications would require substantial effort and may not be justified for a primarily confirmatory study.

The second possibility is a meta-epidemiological study estimating how often network meta-analyses include multicomponent interventions and how often their component structure is analyzed explicitly. This could become a short descriptive paper and a foundation for later methodological research, but it was not adopted as an immediate task.

The agreed priority is to look first for reusable datasets in archived network-meta-analysis databases or other repositories. If no suitable data can be recovered, the reanalysis should be deferred rather than displacing higher-priority methodological work. Bulun will continue reading in CNMA and related methods and will revisit the dissertation direction in a later meeting. Lifeng emphasized that Bulun's longer-term effort should concentrate on statistical method development while using the current AI-audit project as a complementary applied study.

\subsection*{Research Support and Coordination}

AI subscriptions used for research may be reimbursable from grant funds. Bulun was shown how to prepare a non-travel expense report and should retain and organize the receipts. Lifeng will ask Kevin which prior subscription charges remain eligible. If intensive model use continues, a higher-tier subscription may be used for research, subject to the same reimbursement process.

Conference participation and student-paper competitions at major statistics meetings were discussed as professional-development opportunities. These are planning options rather than commitments. Funding was described as generally sufficient for approximately one or two meetings per year, subject to the quality of the work, advance approval, and applicable travel rules.

\section*{Immediate Action Plan}

\begin{enumerate}
    \item \textbf{Bulun --- Bayesian $I^2$ manuscript (target: within one to two weeks).} Merge the new case study into the main manuscript; replace the selected existing example; revise the explanation of case-study selection; integrate the forest plot; reconcile the code with the shared analysis files; and proofread the complete manuscript.

    \item \textbf{Bulun --- TAILOR final audit.} Run additional AI-assisted checks, verify every substantive flag manually, make any warranted corrections, and notify Lifeng when the manuscript is ready.

    \item \textbf{Lifeng --- TAILOR submission.} After Bulun's notification, complete the final review, draft and share the cover letter, and submit the manuscript.

    \item \textbf{Bulun --- publication-audit prompt and corpus.} Remove the fixed ten-finding limit, refine the definitions of eligible errors, preserve a separate uncertain category, and rerun the workflow across the agreed 2018--2019 and 2024--2025 periods.

    \item \textbf{Bulun and Lifeng --- verification meeting.} When the reruns are substantially complete, arrange a meeting with Shiying, preferably within one to two weeks. At that meeting, assign human verification across both time periods and discuss manuscript structure, authorship responsibilities, and possible journals.

    \item \textbf{Bulun --- audit outputs.} Preserve the detailed reasoning and source locations for supplementary material; organize results by year and paper; and identify a concise set of independently verified examples for the main manuscript.

    \item \textbf{Bulun and Lifeng --- AI subscription reimbursement.} Bulun will collect recent receipts and complete the expense-report entries. Lifeng will ask Kevin about the allowable retrospective reimbursement period.
\end{enumerate}

\section*{Future Work}

\begin{itemize}
    \item Convert the Bayesian $I^2$ procedure into a documented R function with reproducible example datasets, then integrate it into Lifeng's existing meta-analysis package after the manuscript and planned package update are complete. Function naming and any broader S3 interface remain open.

    \item Develop the publication-audit report into a paper with a transparent verification protocol, balanced time periods, article-level supplementary records, and carefully qualified temporal interpretation. The target journal and final division of labor remain to be decided.

    \item Determine whether the strongest audit findings have consequences beyond local presentation. Any severity assessment should be evidence-based and should avoid overstating the effect of minor reporting discrepancies.

    \item Search for an analysis-ready multicomponent network-meta-analysis dataset. If one is available, compare traditional NMA and CNMA; otherwise defer manual extraction or assign data collection as a future project.

    \item Consider a meta-epidemiological study of the prevalence and handling of multicomponent interventions in published network meta-analyses as a possible short paper and methodological foundation.

    \item Continue reviewing CNMA and related methodology, identify a dissertation direction centered on statistical development, and revisit the choice in a subsequent advising meeting.

    \item Monitor suitable student-paper competitions and conference opportunities once a sufficiently mature project is available; no conference submission or travel plan was finalized at this meeting.
\end{itemize}

\end{document}
